\section{Discussion}

\subsection{Why a real road network was necessary}

Each of our main findings is invisible on the graph families that dominate
this literature, and for a structural reason rather than a matter of scale.

The maximum-flux walk fails by \emph{dead-ending}: it reaches a node with no
admissible unvisited successor. On a complete graph that is impossible, since
every unvisited node is adjacent; on a small-world graph of mean degree 6 to
600 it is vanishingly rare. Our network has mean out-degree 2.46, and the
undirected pressure solve routinely prefers a direction the road forbids. The
failure mode we spent most of this study measuring is structurally
unreachable on the substrates where the model's accuracy was established.

This reframes the disagreement with prior work as a question of domain rather
than of correctness. Zhang et al.\ report 100\% accuracy with \emph{no}
extraction machinery at all~\cite{zhang2014improved}, and they are entitled
to: on Erd\H{o}s--R\'enyi graphs the flux separates cleanly and reading the
path off is trivial. Wang et al.\ state that the dominant path at
convergence \emph{is} the shortest path~\cite{wang2021transition}, on
explicitly undirected graphs where the walk cannot dead-end. Neither result
is wrong. What changes on a directed, sparse, organically grown network is
that extraction stops being a formality and becomes the step that determines
the reported outcome.

\subsection{Provenance as a reporting requirement}

The practical consequence is a reporting rule. A Physarum accuracy figure on
a network of this kind is a statement about the extraction policy at least as
much as about the model, and a pooled number is unfalsifiable: our own
pooled 90\% and model-only 23.3\% come from the same solves. We would argue
that any Physarum evaluation on a directed network should publish the
provenance distribution of its reported paths, and that a model whose
fallback is a shortest-path algorithm must report what share of its results
that fallback produced.

The warm-start case shows how far this can go. At 66.8\% full-graph fallback,
a warm Physarum recovery is more likely to be Dijkstra's answer than the
model's, and an evaluation reporting only pooled stretch would record that
configuration as near-perfect. It was precisely such a configuration --- the
stranded-component degeneracy of Section~\ref{sec:algorithms} --- that
produced a fabricated mean stretch of exactly 1.000 on the scenario where
the model is weakest, with no error raised. Instrumenting provenance is what
made that visible.

\subsection{What self-healing is actually worth here}

Our answer to the motivating question is a quantified trade rather than an
endorsement or a dismissal.

The adaptive state is real and it is substantial: warm start recovers in a
single iteration on nearly half of all closures, which cold start never
achieves, and per-iteration cost is unchanged, so the saving is genuinely
adaptation. That is the strongest form of the self-healing claim we can
support, and it is stronger than any ratio we could quote.

But it is paid for. Warm start is significantly worse on the reported
post-disruption path, and it degrades the model's own extraction further
still. A system that must guarantee route quality after a closure should
prefer a cold solve or a recompute; one that must respond within a single
adaptation step, and can tolerate a few percent of path cost, is the case
where retained state pays. Neither conclusion is available from the speed
measurement alone.

Against all of this stands the baseline. Dijkstra recomputes an exact route
in about 31\,ms on this network, and none of the bio-inspired methods comes
close on wall-clock. For a graph of a few thousand nodes, recomputation is
simply the right engineering answer, and the interest in these methods has to
rest on regimes we did not test --- far larger networks, many simultaneous
commodities, or settings where the cost function is not globally observable.
We think it is more useful to say so than to construct a comparison the
baseline would lose.

\subsection{A methodological result about benchmarking}

The congestion-model finding generalises beyond this study. We measured
opposite signs for the effect of congestion on ACO from an analytic
volume-delay function and from a microsimulator, on the same network with the
same OD pairs. Both are standard and defensible choices; they differ because
SUMO puts signal delay in the numerator but not the free-flow denominator,
while BPR loads synthetic demand onto shortest paths, which is exactly where
an OD-path measurement walks.

A benchmark that reports results ``under congestion'' without naming the
source is therefore not reproducible in the sense that matters, because an
equally reasonable alternative can reverse the sign of its headline effect.
We would encourage reporting both where feasible, treating them as brackets
rather than as substitutes, and never merging them into a single congestion
level.
